%! AP Statistics — Unit 1, Lesson 1.1 — Guided Notes ANSWER KEY
\documentclass[10pt]{article}
\usepackage[margin=0.65in, top=0.55in, bottom=0.65in]{geometry}
\usepackage{xcolor}
\usepackage[most]{tcolorbox}
\usepackage{enumitem}
\usepackage{multicol}
\usepackage{amsmath,amssymb}
\usepackage{array}
\usepackage{tabularx}
\usepackage{tikz}
\usepackage{colortbl}
\usepackage{amssymb}

\tcbuselibrary{breakable, skins}

% ── Colors ──────────────────────────────────────────────────────────────────
\definecolor{navy}{HTML}{1F3A5F}
\definecolor{navylight}{HTML}{2E5490}
\definecolor{sky}{HTML}{EAF3FB}
\definecolor{skymid}{HTML}{C5DFF5}
\definecolor{goldacc}{HTML}{D4820A}
\definecolor{hookbg}{HTML}{FEF3E2}
\definecolor{greenbg}{HTML}{EDF8F0}
\definecolor{greenacc}{HTML}{2D7A4F}
\definecolor{redbg}{HTML}{FCECEC}
\definecolor{redacc}{HTML}{C0392B}
\definecolor{charcoal}{HTML}{2C3E50}
\definecolor{slate}{HTML}{64748B}
\definecolor{linegray}{HTML}{AAAAAA}
\definecolor{keyred}{HTML}{CC0000}

% ── Box styles ───────────────────────────────────────────────────────────────
\newtcolorbox{objectivebox}{
    colback=sky, colframe=navy, boxrule=1pt, arc=2mm,
    left=3mm, right=3mm, top=2mm, bottom=2mm,
    fonttitle=\bfseries\small\color{white}, title={Primary Objective},
    attach boxed title to top left={yshift=-2mm, xshift=4mm},
    boxed title style={colback=navy, colframe=navy, arc=1.5mm}
}
\newtcolorbox{vocabbox}{
    colback=greenbg, colframe=greenacc, boxrule=0.8pt, arc=2mm,
    left=3mm, right=3mm, top=2mm, bottom=2mm,
    fonttitle=\bfseries\small\color{white}, title={Vocabulary \& Key Concepts},
    attach boxed title to top left={yshift=-2mm, xshift=4mm},
    boxed title style={colback=greenacc, colframe=greenacc, arc=1.5mm},
    breakable
}
\newtcolorbox{hookbox}{
    colback=hookbg, colframe=goldacc, boxrule=0.8pt, arc=2mm,
    left=3mm, right=3mm, top=2mm, bottom=2mm,
    fonttitle=\bfseries\small\color{white}, title={Hook},
    attach boxed title to top left={yshift=-2mm, xshift=4mm},
    boxed title style={colback=goldacc, colframe=goldacc, arc=1.5mm},
    breakable
}
\newtcolorbox{notesbox}[1]{
    colback=white, colframe=navy, boxrule=0.8pt, arc=2mm,
    left=3mm, right=3mm, top=2mm, bottom=2mm,
    fonttitle=\bfseries\small\color{white}, title={#1},
    attach boxed title to top left={yshift=-2mm, xshift=4mm},
    boxed title style={colback=navy, colframe=navy, arc=1.5mm},
    breakable
}
\newtcolorbox{practicebox}{
    colback=redbg, colframe=redacc, boxrule=0.8pt, arc=2mm,
    left=3mm, right=3mm, top=2mm, bottom=2mm,
    fonttitle=\bfseries\small\color{white}, title={Guided Practice},
    attach boxed title to top left={yshift=-2mm, xshift=4mm},
    boxed title style={colback=redacc, colframe=redacc, arc=1.5mm},
    breakable
}
\newtcolorbox{spiralbox}{
    colback=white, colframe=slate, boxrule=0.6pt, arc=2mm,
    left=3mm, right=3mm, top=2mm, bottom=2mm,
    fonttitle=\bfseries\small\color{white}, title={Connections \& Big Ideas},
    attach boxed title to top left={yshift=-2mm, xshift=4mm},
    boxed title style={colback=slate, colframe=slate, arc=1.5mm}
}

% ── Answer helpers ───────────────────────────────────────────────────────────
% \ans{text} — inline answer replacing a \blank
\newcommand{\ans}[1]{\underline{\textcolor{keyred}{\textbf{#1}}}}

% \ansline{text} — answer replacing a write line
\newcommand{\ansline}[1]{%
  \vspace{1pt}\noindent\textcolor{keyred}{\textbf{#1}}%
  \textcolor{linegray}{\dotfill}\vspace{3pt}}

% \writeline — blank line (kept for structural spacing)
\newcommand{\writeline}{\vspace{1pt}\noindent\textcolor{linegray}{\hrulefill}\vspace{3pt}}

% ── List / spacing settings ──────────────────────────────────────────────────
\setlist[itemize]{leftmargin=1.2em, itemsep=0.3em, topsep=0.3em}
\setlist[enumerate]{leftmargin=1.5em, itemsep=0.4em, topsep=0.3em}
\setlength{\parindent}{0pt}
\setlength{\parskip}{3pt}

%=============================================================================
\begin{document}

% ── Page header ──────────────────────────────────────────────────────────────
\begin{tikzpicture}[remember picture, overlay]
  \fill[navy] (current page.north west)
    rectangle ([yshift=-0.52in]current page.north east);
\end{tikzpicture}
\vspace{-0.42in}
{\color{white}\textbf{\large AP Statistics \hfill Unit 1: Exploring One-Variable Data}}\\[2pt]
{\color{skymid}\small\textbf{Lesson 1.1 \quad Introducing Statistics: What Can We Learn from Data?}}
\hfill
{\color{keyred}\small\textbf{--- ANSWER KEY ---}}
\vspace{0.18in}

% Name / Date / Period row
\noindent
\begin{tabularx}{\linewidth}{X X X}
Name:\;\ans{\;Key\;} & Date:\;\ans{\;\textit{(any)}\;} & Period:\;\ans{\;\textit{(any)}\;}
\end{tabularx}

\vspace{0.12in}

% ── OBJECTIVE ────────────────────────────────────────────────────────────────
\begin{objectivebox}
\small By the end of this lesson, I will be able to\dots\\[4pt]
\ansline{identify individuals, populations, samples, and variables in a statistical context.}\\
\ansline{explain why statistics is needed — because individuals vary, and data captures that variation.}\\
\ansline{describe the four-step statistical problem-solving process.}
\end{objectivebox}

\vspace{0.1in}

% ── VOCABULARY ───────────────────────────────────────────────────────────────
\begin{vocabbox}
\small
\begin{multicols}{2}

\noindent\textbf{\textcolor{navy}{Statistics:}}\quad
\ans{the science of collecting, analyzing,}\\
\ans{and drawing conclusions from data}\\[4pt]

\noindent\textbf{\textcolor{navy}{Individual:}}\quad
\ans{a single person, object, or event}\\
\ans{in a dataset}\\[4pt]

\noindent\textbf{\textcolor{navy}{Population:}}\quad
\ans{the entire group of individuals}\\
\ans{of interest}\\[4pt]

\noindent\textbf{\textcolor{navy}{Sample:}}\quad
\ans{a subset of the population}\\
\ans{from which data are collected}\\[4pt]

\noindent\textbf{\textcolor{navy}{Census:}}\quad
\ans{data collected from EVERY individual}\\
\ans{in the population}\\[4pt]

\columnbreak

\noindent\textbf{\textcolor{navy}{Variability / Variation:}}\\
\textcolor{keyred}{\textbf{Natural differences in values across}}\\
\textcolor{keyred}{\textbf{individuals — the reason statistics exists}}\\[4pt]

\noindent\textbf{\textcolor{navy}{Data:}}\quad
\ans{recorded values for individuals;}\\
\ans{can be categorical or quantitative}\\[4pt]

\noindent\textbf{\textcolor{navy}{Anecdotal evidence:}}\\
\textcolor{keyred}{\textbf{data from a few unrepresentative cases;}}\\
\textcolor{keyred}{\textbf{not reliable for drawing conclusions}}\\[4pt]

\end{multicols}
\end{vocabbox}

\vspace{0.1in}

% ── HOOK ─────────────────────────────────────────────────────────────────────
\begin{hookbox}
\small\textbf{Headline:}\quad
\textit{``Students who sleep more score higher on tests.''}
\hfill\textcolor{slate}{\small(national education study, $n = 12{,}000$)}

\vspace{6pt}
\textbf{Discussion questions — write your thoughts:}
\begin{enumerate}[leftmargin=1.5em, itemsep=4pt]
  \item How do researchers know this? Did they study every student in the world?\\
        \textcolor{keyred}{\textbf{No — they used a sample of 12,000 students and applied statistical}}\\
        \textcolor{keyred}{\textbf{methods to draw conclusions about a larger population.}}

  \item Could this finding be wrong? What would make you doubt it?\\
        \textcolor{keyred}{\textbf{Yes — the sample might not be representative; other factors (e.g., overall}}\\
        \textcolor{keyred}{\textbf{health, study habits) could explain both sleep and test performance.}}

  \item What information would you want before trusting this headline?\\
        \textcolor{keyred}{\textbf{Who was sampled, how sleep was measured, whether other variables}}\\
        \textcolor{keyred}{\textbf{(confounders) were controlled for, whether the study was an experiment.}}
\end{enumerate}
\end{hookbox}

\vspace{0.1in}

% ── DIRECT INSTRUCTION: 4-STEP PROCESS ───────────────────────────────────────
\begin{notesbox}{The Statistical Problem-Solving Process}
\small Fill in each step and what it involves.
\vspace{4pt}

\renewcommand{\arraystretch}{1.5}
\begin{tabularx}{\linewidth}{@{} >{\centering\arraybackslash\columncolor{navy}}p{0.18\linewidth} @{\quad$\longrightarrow$\quad} X @{}}
\textcolor{white}{\textbf{Step 1}} &
  \textbf{Ask a \ans{question}} \quad
  \textcolor{slate}{\scriptsize Pose a clear question about a \ans{population} of interest.}\\

\textcolor{white}{\textbf{Step 2}} &
  \textbf{Collect \ans{data}} \quad
  \textcolor{slate}{\scriptsize Select a \ans{sample} and measure \ans{variables} on each individual.}\\

\textcolor{white}{\textbf{Step 3}} &
  \textbf{Analyze the \ans{data}} \quad
  \textcolor{slate}{\scriptsize \ans{Summarize} and \ans{display} what you found.}\\

\textcolor{white}{\textbf{Step 4}} &
  \textbf{Interpret \ans{results}} \quad
  \textcolor{slate}{\scriptsize Draw \ans{conclusions}; acknowledge what the data \ans{cannot tell us}.}\\
\end{tabularx}
\end{notesbox}

\vspace{0.1in}

% ── CONCEPT NOTES ─────────────────────────────────────────────────────────────
\begin{notesbox}{Key Concepts}
\small
\begin{multicols}{2}

\textbf{\textcolor{navy}{Population vs. Sample}}\\
\textit{Example:} A researcher surveys 200 adults about exercise habits.
\begin{itemize}[itemsep=2pt]
  \item Individual: \ans{an adult}
  \item Population: \ans{all adults (e.g., in the U.S.)}
  \item Sample: \ans{the 200 randomly chosen adults}
  \item Variable: \ans{exercise habits / days per week}
\end{itemize}

\vspace{4pt}
\textbf{Why use a sample instead of a census?}\\
\textcolor{keyred}{\textbf{A census is too costly, time-consuming, or impossible for large populations.
A well-chosen sample can give reliable estimates.}}

\columnbreak

\textbf{\textcolor{navy}{Why Statistics Needs Variation}}\\[3pt]
Complete the statement:\\[4pt]
\textit{``If every individual were identical, we would \ans{NOT}
need statistics, because there would be \ans{nothing}
\ans{different} to explain or measure.''}

\vspace{6pt}
\textbf{What does ``variation'' look like in data?}\\
\textit{Example from class — sleep minutes:}\\[4pt]
\begin{center}
{\color{keyred}\textbf{\textit{(record actual class data here)}}}\\[2pt]
\begin{tabular}{|c|c|c|c|c|}
\hline
\textcolor{keyred}{\scriptsize varies} & \textcolor{keyred}{\scriptsize varies} &
\textcolor{keyred}{\scriptsize varies} & \textcolor{keyred}{\scriptsize varies} &
\textcolor{keyred}{\scriptsize varies} \\
\hline
\end{tabular}
\end{center}
\vspace{2pt}
\textit{Min:\;\ans{low}\quad Max:\;\ans{high}\quad ``Typical'':\;\ans{middle value}}

\end{multicols}
\end{notesbox}

\vspace{0.1in}

% ── GUIDED PRACTICE ──────────────────────────────────────────────────────────
\begin{practicebox}
\small\textbf{Scenario Analysis} \hfill \textcolor{slate}{\scriptsize Sample: Hospital knee-surgery study}

\vspace{5pt}
\textit{Scenario:} \textcolor{keyred}{\textbf{A hospital tracks recovery times for 50 patients after knee surgery.}}

\vspace{4pt}
\begin{tabularx}{\linewidth}{@{} l X @{}}
\textbf{Individual:}   & \textcolor{keyred}{\textbf{a single patient who had knee surgery}} \\[6pt]
\textbf{Population:}   & \textcolor{keyred}{\textbf{all patients who have knee surgery (at this hospital / in general)}} \\[6pt]
\textbf{Sample:}       & \textcolor{keyred}{\textbf{the 50 patients whose recovery times were tracked}} \\[6pt]
\textbf{Variable(s):}  & \textcolor{keyred}{\textbf{recovery time (quantitative); could also include pain level, age, etc.}} \\[6pt]
\end{tabularx}

\vspace{4pt}
\textbf{Is a census feasible? Why or why not?}\\
\textcolor{keyred}{\textbf{No — tracking every knee-surgery patient would be prohibitively expensive and}}\\
\textcolor{keyred}{\textbf{time-consuming. The 50-patient sample can still yield useful information.}}

\vspace{4pt}
\textbf{Record one key insight from another group's scenario:}\\
\textcolor{keyred}{\textbf{(Answers will vary — students should capture a specific example of how}}\\
\textcolor{keyred}{\textbf{identifying the population vs. sample changes how we interpret results.)}}
\end{practicebox}

\vspace{0.1in}

% ── SPIRAL / BIG IDEAS ────────────────────────────────────────────────────────
\begin{spiralbox}
\small
\begin{multicols}{2}

\textbf{This lesson connects forward to\dots}
\begin{itemize}[itemsep=2pt]
  \item \textbf{Lesson 1.2:} Types of variables (categorical vs.\ quantitative)
  \item \textbf{Units 6--9:} Using samples to make \emph{inferences} about populations
  \item \textbf{Unit 3:} How we \emph{collect} samples matters for conclusions
\end{itemize}

\columnbreak

\textbf{Big Idea — Write it in your own words:}\\
\textit{What is the point of statistics?}\\[4pt]
\textcolor{keyred}{\textbf{Statistics gives us tools to learn about a large}}\\
\textcolor{keyred}{\textbf{population using a smaller sample, and to}}\\
\textcolor{keyred}{\textbf{account for the natural variation in data.}}

\end{multicols}
\end{spiralbox}

\vspace{0.1in}

% ── EXTRA NOTES ──────────────────────────────────────────────────────────────
\begin{notesbox}{Additional Notes}
\vspace{2pt}
\textcolor{keyred}{\small\textit{Key reminders for grading / discussion:}}\\[2pt]
\textcolor{keyred}{\small $\bullet$\ \textbf{Association $\neq$ causation} — sleeping more may not directly cause higher scores; other variables may be involved.}\\[2pt]
\textcolor{keyred}{\small $\bullet$\ A \textbf{sample} is a group of individuals; \textbf{data} are the values recorded — these are different things.}\\[2pt]
\textcolor{keyred}{\small $\bullet$\ Variation is \textit{expected} and \textit{informative} — it is not a flaw in the data collection.}\\[2pt]
\textcolor{keyred}{\small $\bullet$\ The 4-step process will spiral throughout the entire course; students should internalize it now.}
\end{notesbox}

\end{document}
